Při zvyšování frekvence u jednoho relaxačního mechanismu:
\begin{itemize}
\item \textbf{Relativní permitivita $\varepsilon'$} klesá ze statické hodnoty $\varepsilon_s$ k vysokofrekvenční mezi $\varepsilon_\infty$, protože dipóly přestávají sledovat změny pole.
\item \textbf{Ztrátové číslo $\varepsilon''$} roste do maxima při mezní frekvenci $f_m$ a potom klesá.
\item \textbf{Ztrátový činitel $\tan\delta=\varepsilon''/\varepsilon'$} má rovněž relaxační maximum a při velmi nízkých i vysokých frekvencích klesá.
\end{itemize}
Mezní frekvenci vyznačuje přerušovaná čára v grafu; s relaxační dobou souvisí vztahem
\[\boxed{f_m=\frac{1}{2\pi\tau}},\qquad\omega_m\tau=1.\]
